\section{Introduction}

Many biological phenomena are observed as development, morphogenesis, differentiation, and evolution. These phenomena are often described using language such as design, program, or genetic instruction. Such language can be useful at an intuitive level, but its relation to the underlying physical and chemical processes is not always explicit.

DNA sequences constrain molecular production and reaction pathways. They do not directly specify a completed structure, purpose, or intention. The fact that cells with the same DNA can differentiate into distinct states illustrates this point. DNA is therefore treated here as a parameter set \(\theta\) that constrains the dynamics of the system.

This paper describes biological systems as processes of state change and sustained structure. We use three quantities: information density \(\I\), information flow \(\J\), and induced potential \(\PhiI\). These quantities provide a common description of development and morphogenesis.

The central feature of the framework is that it does not require design or purpose as explanatory primitives. Form and function appear as outcomes of state changes that converge toward stable structures.

The aim is to describe existing phenomena in a consistent form. We combine a minimal set of variables, a continuity equation, and an attractor-based interpretation of stability.

Section 2 gives the intuitive picture. Section 3 defines the formal mapping. Section 4 introduces the core equations. Sections 5 and following apply the framework to biological phenomena, simulation, and scope.

This paper builds on the Information-Field-Geometry-Theory (IFGT), which formulates the motion of systems toward closure and the persistence of quasi-closure structures. In IFGT, the degree of closure \(C\) characterizes how far a system has progressed toward a stable configuration. In biological systems, the field \(\I\) should not be identified directly with the IFGT quantity \(I=f(1-C)\). It is a biological, coarse-grained constraint-density field: a measure of local structural constraint, differentiation, and retained bias within a particular observational window. The relation is therefore a correspondence between descriptive layers, not an identity between variables. The sustained non-equilibrium structure of living systems corresponds, in IFGT terms, to quasi-closure: a regime that remains on the generated side of the complete-closure boundary and persists because residual difference continues to be available for update.
